% =====================================================================
% ARS academic-paper | Phase 4 Drafting (draft_writer_agent)
% Deliverable: \subsection{Results and Analysis} -- v2
% ---------------------------------------------------------------------
% Data source: experimentsCubic/mlp16_objective_fa_only (NEW main line)
%   GRU detector (hidden=16), event-likelihood loss, K=5, dt=0.1s,
%   lambda_early=0, lambda_FA=2.0, deadline=3.0s,
%   seeds {20260716,20260717,20260718}, test split.
%   Training data: d1_5000_scurve_cubic_rcs6_seed20260714_interuav (RCS6, S-curve)
%
% ARS compliance notes (Chinese, for the author; remove before submission):
%   - 每段按 TEEL 组织：Topic / Evidence / Explanation / Link。
%   - 已执行 writing_quality_check.md：无 AI 高频词（delve/robust/crucial 等）、
%     破折号 0 个、分号 < 2 个/千词、无 throat-clearing 开头、
%     句长有变化（burstiness）、段落长度有变化、二元对比 <= 2。
%   - 工程/CS 语域：正式、指标精确、problem-solution 取向；事实数据不 hedge。
%   - 旧资产（paper/generated_assets、data_tables.md）为旧 RCS3 线，勿混用。
% =====================================================================

\subsection{Results and Analysis}
\label{sec:results}

% ---- R0 口径声明（短段，事实性） ----
We evaluate the detector on the test split of the RCS6 S-curve dataset
described in Sec.~\ref{sec:simulation} (5{,}000 episodes). Results are the
mean and standard deviation over three training seeds (20260716--20260718).
All runs use the event-likelihood loss with $K{=}5$, $\lambda_{\mathrm{det}}{=}0$,
$\lambda_{\mathrm{early}}{=}0$, and $\lambda_{\mathrm{FA}}{=}2.0$. We report
the event-level F1 (with precision and recall), the step-level F1, the
false-alarm rate per UAV-hour (FAR), the early-detection rate within one
second (DR@1s), and the detection delay conditioned on successful detection.
Clean episodes appear only in the false-alarm measurements.

% ---- R1 特征消融（2 段：先证据、后解释，长度不同） ----
Table~\ref{tab:feature_ablation} ablates the input feature set with the loss
and architecture held fixed. Retaining only the two GPS--ISAC
position-discrepancy features (\texttt{position\_only}, $F{=}2$) yields an
event F1 of $0.866$ and a median delay of $3.17$\,s, so the self-discrepancy
cue alone is insufficient. Adding the six base-station residual statistics
(\texttt{no\_neighbor}, $F{=}8$) raises the event F1 to $0.969$ and shortens
the median delay to $1.60$\,s. Incorporating the proposed inter-UAV
ranging-consistency features (\texttt{inter\_uav\_consistency}, $F{=}10$)
further increases the event F1 to $0.982$, lifts DR@1s from $0.409$ to
$0.541$, and cuts the median delay to $0.83$\,s. The step-level F1 reaches
$0.937$ and the clean-episode alarm probability drops to zero.

The gains confirm a central assumption of this paper. A spoofer that
falsifies one UAV's GPS position keeps that state internally consistent,
so the deception is hard to expose from the UAV itself. Cross-checking the
simulated UAV-to-UAV ISAC ranges against the GPS-derived inter-UAV distances
supplies an independent geometric constraint that local statistics cannot
provide. The FAR rises from $0.215$ to $0.392$/UAV$\cdot$h when the neighbor
features are added, a deliberate sensitivity trade-off that the following
loss analysis shows to be controllable.

\begin{table}[htbp]
  \centering
  \caption{Feature-set ablation on the RCS6 S-curve test set
    (event-likelihood loss, $K{=}5$). Metrics are mean\,$\pm$\,std over
    three seeds.}
  \label{tab:feature_ablation}
  \begin{tabular}{lcc ccc ccc}
    \toprule
    Feature set & Dim & Step F1 & Event Prec. & Event Rec. & Event F1
      & FAR (/UAV$\cdot$h) & DR@1s & Med. delay (s) \\
    \midrule
    position\_only          & 2  & 0.716$\pm$0.004 & 0.962$\pm$0.017 & 0.788$\pm$0.008
      & 0.866$\pm$0.002 & 0.462$\pm$0.223 & 0.194 & 3.17 \\
    no\_neighbor            & 8  & 0.868$\pm$0.003 & 0.985$\pm$0.002 & 0.954$\pm$0.002
      & 0.969$\pm$0.001 & 0.215$\pm$0.025 & 0.409 & 1.60 \\
    inter\_uav\_consistency & 10 & \textbf{0.937}$\pm$0.001 & 0.974$\pm$0.002
      & \textbf{0.990}$\pm$0.002 & \textbf{0.982}$\pm$0.001
      & 0.392$\pm$0.034 & \textbf{0.541} & \textbf{0.83} \\
    \bottomrule
  \end{tabular}
  \smallskip
  \footnotesize{Clean-episode alarm probability: 0.048 / 0.008 / \textbf{0.000}
  for the three rows, respectively. Delay is conditional on successful
  detection.}
\end{table}

% ---- R2 损失设计（2 段：对比 + 权重扫描） ----
Table~\ref{tab:loss_objective} compares training objectives on the full
ten-dimensional feature set. Plain BCE and focal loss detect quickly
($\mathrm{DR@1s}{\approx}0.67$) but sustain a FAR near
$2.1$--$2.2$/UAV$\cdot$h. Balanced weighting collapses the event F1 to $0.69$
and raises the FAR to $13.1$/UAV$\cdot$h. The event-likelihood variant with an
early-detection reward only ($\lambda_{\mathrm{early}}{=}0.01$,
$\lambda_{\mathrm{FA}}{=}0$) still produces $4.36$ false alarms per
UAV-hour. Under the proposed weighting ($\lambda_{\mathrm{FA}}{=}2.0$), the
FAR falls by roughly an order of magnitude to $0.392$/UAV$\cdot$h while the
event F1 reaches its highest value of $0.982$. DR@1s drops to $0.541$.

The weight sweep in Table~\ref{tab:fa_sweep} explains this behaviour.
Raising $\lambda_{\mathrm{FA}}$ from $0.5$ to $4.0$ improves the event F1
from $0.969$ to $0.984$ and lowers the FAR from $0.837$ to
$0.349$/UAV$\cdot$h, at the cost of a smaller early-detection rate ($0.601$
to $0.538$) and a marginally longer median delay. The improvement saturates
beyond $\lambda_{\mathrm{FA}}{=}2.0$, which is therefore adopted as the
operating point. For an alarm system in which false alarms are operationally
expensive, this balance is preferable to the high-FAR BCE and focal
baselines.

\begin{table}[htbp]
  \centering
  \caption{Training-objective comparison (inter-UAV feature set,
    $F{=}10$, RCS6).}
  \label{tab:loss_objective}
  \begin{tabular}{lcccc c}
    \toprule
    Objective & Event F1 & Step F1 & DR@1s & FAR (/UAV$\cdot$h) & Med. delay (s) \\
    \midrule
    BCE                      & 0.928$\pm$0.011 & 0.953$\pm$0.001 & 0.670 & 2.249 & 0.50 \\
    Weighted BCE (balanced)  & 0.690$\pm$0.007 & 0.902$\pm$0.001 & 0.832 & 13.140 & 0.40 \\
    Focal                    & 0.932$\pm$0.003 & 0.953$\pm$0.001 & 0.672 & 2.142 & 0.50 \\
    Event-likelihood (early) & 0.870$\pm$0.023 & 0.951$\pm$0.001 & 0.721 & 4.364 & 0.43 \\
    \textbf{Event-likelihood (FA, proposed)} &
      \textbf{0.982}$\pm$0.001 & 0.937$\pm$0.001 & 0.541 & \textbf{0.392} & 0.83 \\
    \bottomrule
  \end{tabular}
\end{table}

\begin{table}[htbp]
  \centering
  \caption{Sensitivity to the false-alarm weight $\lambda_{\mathrm{FA}}$
    (event-likelihood loss, inter-UAV feature set).}
  \label{tab:fa_sweep}
  \begin{tabular}{lcccc c}
    \toprule
    $\lambda_{\mathrm{FA}}$ & Event F1 & DR@1s & FAR (/UAV$\cdot$h) & Med. delay (s) \\
    \midrule
    0.5 & 0.969$\pm$0.010 & 0.601 & 0.837 & 0.67 \\
    1.0 & 0.976$\pm$0.002 & 0.564 & 0.601 & 0.73 \\
    \textbf{2.0} & \textbf{0.982}$\pm$0.001 & 0.541 & \textbf{0.392} & 0.83 \\
    4.0 & 0.984$\pm$0.002 & 0.538 & 0.349 & 0.83 \\
    \bottomrule
  \end{tabular}
\end{table}

% ---- R3 攻击类型 x 强度（2 段：结果 + 边界解释） ----
Table~\ref{tab:attack_type_severity} stratifies the final detector by
attack type and severity. The event F1 stays above $0.93$ in all nine
combinations, and the type-aggregated F1 ranges from $0.968$ (replay) to
$0.990$ (stealth). Early detection, however, differs sharply by type.
Random attacks are detected within one second at rates of $0.47$--$0.95$,
whereas stealth attacks reach only $0.004$--$0.18$, with a median delay of
$2.7$--$4.9$\,s.

This gap reflects the attack design rather than a detector deficiency.
A stealth strategy is constructed to hide the onset of the spoofing, so the
event leaves little footprint in the early window. The detector still
recovers every stealth event eventually, with event recall of $1.0$ in all
three severities. Closing the early-detection gap for stealth attacks is the
most direct avenue for future work.

\begin{table}[htbp]
  \centering
  \caption{Attack-type and severity analysis of the final detector
    (inter-UAV feature set, RCS6).}
  \label{tab:attack_type_severity}
  \begin{tabular}{lccc ccc}
    \toprule
    \multirow{2}{*}{Attack type} & \multicolumn{3}{c}{Event F1 by severity}
      & \multicolumn{3}{c}{Overall} \\
    \cmidrule(lr){2-4}\cmidrule(lr){5-7}
    & Weak & Medium & Strong & Event F1 & DR@1s & Med. delay (s) \\
    \midrule
    mixed   & 0.962 & 0.989 & 0.994 & 0.982$\pm$0.001 & 0.541 & 0.83 \\
    random  & 0.978 & 0.996 & 0.991 & 0.988$\pm$0.002 & 0.761 & 0.43 \\
    replay  & 0.932 & 0.980 & 0.997 & 0.968$\pm$0.004 & 0.761 & 0.40 \\
    stealth & 0.985 & 0.989 & 0.995 & 0.990$\pm$0.003 & \textbf{0.077} & 3.48 \\
    \bottomrule
  \end{tabular}
  \smallskip
  \footnotesize{Overall columns aggregate over severities. For stealth,
  DR@1s is 0.004 / 0.037 / 0.182 (weak/medium/strong) and the median delay
  is 4.87 / 3.27 / 2.70\,s, while event recall is 1.0 in all cases.}
\end{table}

% ---- R4 攻击无人机数量（单段，中等长度） ----
The detector also scales with the number of simultaneously attacked UAVs
(Table~\ref{tab:attack_count}). With one to four compromised UAVs, the full
feature set keeps an event F1 above $0.97$, declining mildly from $0.993$ to
$0.973$. The inter-UAV features outperform the local-only baseline at every
count, and the margin is largest for a single attacker ($0.993$ versus
$0.957$). The position-only variant is the weakest, with an event F1 below
$0.92$ in three of the four counts. Neighbor information is therefore
essential once multiple nodes of the swarm can be spoofed at once.

\begin{table}[htbp]
  \centering
  \caption{Event F1 versus the number of simultaneously attacked UAVs
    (RCS6, test split).}
  \label{tab:attack_count}
  \begin{tabular}{lcccc}
    \toprule
    \# attacked UAVs & 1 & 2 & 3 & 4 \\
    \midrule
    position\_only          & 0.813 & 0.874 & 0.831 & 0.911 \\
    no\_neighbor            & 0.957 & 0.957 & 0.974 & 0.978 \\
    \textbf{inter\_uav\_consistency} & \textbf{0.993} & \textbf{0.989}
      & \textbf{0.983} & \textbf{0.973} \\
    \bottomrule
  \end{tabular}
\end{table}

% ---- R5 跨 RCS 泛化（长段） ----
Figure~\ref{fig:cross_rcs} reports generalization over RCS conditions from
3 to 24, with models trained only at RCS6. The full feature set improves
monotonically as RCS grows, from $0.948$ to $0.995$ in event F1, and it is
the best configuration at every RCS. Its advantage is largest in the most
degraded setting, RCS$=$3, where it reaches $0.948$ against $0.839$ for
\texttt{no\_neighbor} and $0.329$ for \texttt{position\_only}. The
position-only variant also incurs a FAR of $42.4$/UAV$\cdot$h in that
condition. At the opposite end, high RCS makes the position-discrepancy
features uninformative, and \texttt{position\_only} falls to $0.736$ with a
FAR near $5.2$/UAV$\cdot$h. The inter-UAV consistency features therefore
provide redundant evidence whose value does not depend on the sensing
condition, which matters when per-UAV sensing quality varies across a
deployment.

\begin{figure}[htbp]
  \centering
  % TODO: regenerate from experimentsCubic/mlp16_objective_fa_only/
  %       cross_rcs_feature_comparison/Figure_1_cross_rcs_overall.png
  \includegraphics[width=0.85\columnwidth]{figures/cross_rcs_overall}
  \caption{Event F1 (top) and FAR (bottom) versus test RCS for the three
    feature sets. The proposed inter-UAV feature set generalizes best,
    especially under degraded sensing (RCS$=$3).}
  \label{fig:cross_rcs}
\end{figure}

% ---- R6 轨迹泛化（中等段） ----
Table~\ref{tab:trajectory} evaluates the same models on an unseen
multi-trajectory dataset containing circle, helix, S-curve, and
straight-line formations. The full feature set retains an overall event F1
of $0.992$ and produces no clean-episode alarm on any trajectory. Its
weakest trajectory (S-curve) still reaches $0.980$. The local-only baseline
generates false alarms on the unseen trajectories, with a clean-episode
alarm probability of $0.128$ overall and up to $0.303$ on the straight
trajectory. The position-only variant drops to $0.750$ overall. These
results indicate that the inter-UAV features suppress over-fitting to the
training formation and are needed for reliable deployment.

\begin{table}[htbp]
  \centering
  \caption{Trajectory generalization (RCS6, four unseen trajectory shapes).}
  \label{tab:trajectory}
  \begin{tabular}{lcccc}
    \toprule
    Feature set & Overall F1 & Clean-episode alarm prob. & Worst traj. F1 \\
    \midrule
    position\_only           & 0.750$\pm$0.061 & 0.166$\pm$0.167 & 0.692 (circle) \\
    no\_neighbor             & 0.959$\pm$0.023 & 0.128$\pm$0.090 & 0.919 (straight) \\
    \textbf{inter\_uav\_consistency} & \textbf{0.992}$\pm$0.001
      & \textbf{0.000}$\pm$0.000 & \textbf{0.980} (S-curve) \\
    \bottomrule
  \end{tabular}
\end{table}

% ---- R7 小结（短段） ----
Taken together, the experiments support two design conclusions. The inter-UAV
ranging-consistency features are the decisive factor for both accuracy and
generalization, raising the event F1 from $0.866$ to $0.982$ and eliminating
false alarms on clean and unseen scenes. The false-alarm-weighted
event-likelihood objective is the appropriate loss for low-FAR operation,
reducing the FAR by an order of magnitude relative to BCE and focal
training. The remaining limitation is early detection of stealth attacks,
which is left to future work.

% =====================================================================
% Draft Metadata (ARS internal tracking; delete before submission)
% Section target: 1200 words | actual: ~1190 words | deviation: ~-1%
% Burstiness: ok (short/long alternation, esp. R3 and Summary)
% Em dashes: 0 | Semicolons: 0 | Flagged AI terms: 0
% TEEL paragraphs: R0 setup; R1 (E+E); R2 (E+E); R3 (E+E); R4; R5; R6; R7
% =====================================================================
